\section{Conclusion}
\label{sec:conclusion}

Ce projet a porté sur Apache Spark l'algorithme StreamKM++, dont une
implémentation de référence en Apache Flink est étudiée par
\cite{bitsakisection section section section 8}. Nous avons implémenté et testé l'intégralité de la
chaîne : le noyau algorithmique non-distribué (seeding $Dsection 2$ pondéré,
algorithme de Lloyd, coreset tree, merge-and-reduce), sa distribution sur
Spark via une fusion en arbre plutôt que séquentielle, et un instrument de
mesure de la qualité du clustering réutilisable pour l'ensemble des
expériences du projet.

Les résultats obtenus sur le jeu de données complet Covertype
section section section 1\section section section 2 points) montrent que notre implémentation distribuée reproduit la
qualité de la référence séquentielle à moins de section 1\,\%$ près, et reste dans
un ordre de grandeur comparable à Spark MLlib tout en s'appuyant sur un
résumé (coreset) de taille très inférieure au jeu de données complet. Une
relecture ciblée du code a permis de corriger un bug réel affectant
l'indépendance de l'aléa dans la fusion distribuée, de dériver rigoureusement
une tolérance de test plutôt que de la choisir arbitrairement, et de combler
un angle mort de couverture de test sur l'état maintenu entre micro-batches.

Deux volets restent à compléter avant la version finale de ce rapport,
signalés explicitement en \S\ref{sec:experiments} plutôt que comblés par des
valeurs inventées : le passage à l'échelle complète de l'expérience de
trade-off qualité/taille de coreset (actuellement disponible seulement sur
un échantillon préliminaire de 2\section section section 0 points), et l'analyse quantitative de
la scalabilité en temps d'exécution (micro-benchmarks JMH et débit Spark),
en cours de réalisation par un autre membre du groupe.
